\scanpage{115}
% line head | text | 64
\mspage{64}
% line 1–2 | text | et la donnée d'un tel foncteur
et \inserted{la donnée d'}un tel foncteur $\psi_0$ revient bien à celle d'un
$H$. Donc

% line 3–13 | text | Corollaire. La catégorie des données
\uline{Corollaire}\quad La catégorie des données de recollement
topossiques $\varphi:\deleted{\unsure{T}}U^\wedge\to T^\wedge$ entre $\operatorname{Top}(\deleted{\unsure{T}}U),\operatorname{Top}(T)$,
telles que \deleted{$\varphi$ soit} le topos recollé soit \inserted{un ``topos \uline{\unsure{paradisiaque}}'', i.e. \unsure{soit}} équivalent à
un topos de la forme $X^\wedge$ (avec $X$ dans $\mathrm{Cat}$),
est équivalente à celle des données de recollement
$H:U^{\mathrm o}\times T\unsure{-|\to}\mathrm{Ens}$ entre $U$ et $T$. \deleted{\unsure{Elle}} La
\deleted{\unsure{pour un}} donnée \inserted{topossique} $\varphi$ associée à $H$ est le foncteur adjoint
à droite du $T^\wedge\longrightarrow U^\wedge$ déduit de $T\xrightarrow{H}U^\wedge$
par prolongement commutant aux $\lim_{\rightarrow}$, ou encore
elle est donnée par
% line 14 | display | φ(F)(t), Hom on U with wedge
\[
\varphi(\underset{\substack{\rotatebox{-90}{$\in$}\\U^\wedge}}{F})(t)
\simeq\operatorname{Hom}_{U^\wedge}(\deleted{\unsure{F},}H(t),F),
\]
% line 15–19 | text | où dans le second membre
où dans le second membre $H$ est interprété comme
un foncteur $T\longrightarrow U^\wedge$. Le topos recollé
\inserted{par $\varphi$} est can. équivalent à $\operatorname{Top}(X)$, où $X$ est
l'objet de $\mathrm{Cat}$ obtenu par recollement à l'aide
de $H$.

% line 20–22 | text | Du point de vue de la théorie des faisceaux
Du point de vue de la théorie des faisceaux,
cela donne l'interprétation de $H$, comme
exprimant le foncteur
% line 23 | display | φ equals j upper star i lower star
\[
\varphi=j^*i_*:U^\wedge-T^\wedge
\]
% line 24 | text | par la formule précédente
par la formule précédente.
% archive sheet number: 115
